\section{The Audited Benchmark}
\label{sec:background}

What does InventoryBench specify, and what did we have to establish ourselves? We
answer from its code and its data, because three answers differ from its
documentation and one of them moves a published number.

\textbf{Inventory and lead times.} There are 1{,}320 instances, 720 synthetic and
600 derived from real weekly apparel sales, split evenly across three lead-time
directories. Synthetic episodes run 50 decisions and real ones 47, not the 48
often reported, which counts the trajectory file's header line. Two lead-time
quantities are easy to conflate: the \emph{promised} lead time, passed to the
policy at initialization with values $\{0, 4, 2\}$, and the \emph{actual}
per-period lead time, never shown to the policy, with values $\{0\}$, $\{4\}$ or
$\{1,2,3,\infty\}$, where $\infty$ encodes a lost shipment. Exactly 440 instances,
the whole stochastic third, contain at least one lost period.

\textbf{Event order and observability.} Each period decides, then receives, then
sells, then charges. Two consequences are load-bearing: an order with actual lead
time zero arrives in the same period, because scheduling precedes the arrival
pop, and the in-transit total shown at $t$ still contains units that land at $t$.
Holding falls on ending inventory and nothing is backordered, which makes
lost-sales units a coherent safety endpoint. The policy sees uncensored previous
demand rather than sales, and in-transit inventory is a scalar with no shipment
ages, so many shipment-level histories are consistent with the same receipts.
That last fact drives a decision in Section~\ref{sec:verify}: the controller may
keep an imputed first-in-first-out ledger, and the verifier may never use it as
evidence.

\textbf{Integer economics, and two absences.} Every demand, profit and holding
cell is an integer, so every sum is exact and our equivalence suite can demand a
maximum absolute difference of exactly $0.0$ rather than a tolerance: a nonzero
difference could only be a semantic disagreement. Holding cost is $1.0$ everywhere
and profit is $19$, $4$ or $1$, so \textbf{the labels \texttt{high},
\texttt{med} and \texttt{low} name the margin, not the holding burden}. Two
absences are ours to fill. There is no upstream order cap, so \collie{} registers
its own $C_t$ and labels it as ours wherever it affects a number, distinct from
the published baseline's internal smoother. And a transit pause is not
expressible in the instance format, because both simulators fix a shipment's
arrival period at order time, so it needs a project-owned sidecar.

\subsection{The harness divergence}
\label{sec:divergence}

The benchmark ships two simulators and they disagree about one thing. In the
legacy game simulator a lost order stays in in-transit inventory forever, since
its arrival period is set to infinity and never popped, while in the submission
path it is never added to in-transit at all. The disagreement never changes the
physics, since replaying a run's orders under the other simulator reproduces the
arrivals, the reward and the lost-unit total exactly. It changes what a policy
\emph{reads}, because a base-stock policy under legacy semantics credits units
that will never arrive, orders less, and loses fewer units.

We declare the submission path authoritative, because it is the contract we are
scored against. Under legacy semantics our reimplementation reproduces the
published $0.4447$ to one unit in the last place and matches all 1{,}320 published
order sequences bit for bit, a far stronger check than a matching aggregate that
many policies could reach. Under the authoritative path the same policy scores
$0.5208$, exactly the 880 deterministic-lead-time instances still match, and all
440 stochastic ones differ. Since 440 is precisely the number of instances the
audit found to contain a lost shipment, the divergence is accounted for with no
residue. The published figure was therefore generated under legacy lost-order
semantics and is not reachable from the submission path, so every table comparing
against it must name the harness.
